\chapter{Read e Write in TCFS}
Nel corso di questo capitolo analizzeremmo in dettaglio il funzionamento delle operazioni di scrittura e lettura all'interno del nostro FS; analizzeremo anche la generazione delle chiavi, i design pattern utilizzati e la logica di base, dando un occhiata al codice di TCFS.

\section{Cenni generali sulla lettura e scrittura dei file in TCFS}
Data la quantità di funzioni crittografice e la complessità di alcune di esse, si è scelto di sviluppare una libreria "ad hoc" per semplificare l'utilizzo delle funzioni di cifratura; la libreria denominata crypt-utils funge a questo scopo, essa espone le seguenti funzioni:

\begin{lstlisting}[style=C, caption=Funzioni di crypt-utils, label=C-crypt-utils-functions.h]
#define ENCRYPT true

#define DECRYPT false

#define IV_SIZE 16 //The fixed size of the initialization vector

extern int do_crypt (int mode, FILE *fp, unsigned char **text, int len, unsigned char *key, unsigned char *iv);

void generate_key (unsigned char *destination);

unsigned char *encrypt_string (unsigned char *plaintext, const char *key, int *encrypted_len);

unsigned char *decrypt_string (unsigned char *base64_ciphertext, const char *key);

int is_valid_key (const unsigned char *key);

const char *encrypt_file_name_with_hex (const char *file, const char *key);

const char *decrypt_file_name_with_hex (const char *enc_file, const char *key);

const char *encrypt_path (const char *path, const char *key);

const char *encrypt_path_and_filename (const char *path, const char *key);

const char *decrypt_path (const char *encrypted_path, const char *key);

const char *decrypt_path_and_filename (const char *encrypted_path, const char *key);
extern unsigned char *generate_iv (void);

\end{lstlisting}
Nel corso delle successive sezioni ci concentreremo sul come i file vengono cifrati, dunque ci interesseremo a do\_crypt e generate\_key.

\subsection{La funzione do\_crypt}
Prima di parlare della cifratura e decifratura dei blocchi è necessario specificare che ambedue le operazioni si ottengono tramite l'utilizzo della funzione do\_crypt, che funge da "facade" per le funzioni interne alla libreria:
\begin{itemize}
    \item \textbf{encrypt\_file\_gcm}: Che si occupa di cifrare il testo in chiaro e scriverlo sul file.
    \item \textbf{decrypt\_file\_gcm}: Che legge il file cifrato e restituisce il testo in chiaro.
\end{itemize}
In particolare do\_crypt opera in base allo stato della variabile \textit{mode} che può assumere valore ENCRYPT o DECRYPT. Questo permette di modificare facilmente le funzioni ci cifratura e decifratura su file nascondendo gran parte della complessità che esse comportano.

\subsection{La cifratura dei blocchi}
La funzione interna encrypt\_file\_gcm si occupa della cifratura dei dati e della scrittura di essi su file. La funzione richiede i seguenti parametri:
    \begin{itemize}
        \item \textbf{FILE *fp}: Il puntatore al file sul quale scrivere, da notare che TCFS scrive sempre e solo in "\textit{append}"
        \item \textbf{unsigned char *plaintext}: La stringa contenente il testo da cifrare.
        \item \textbf{unsigned char *plaintext\_len}: La lunghezza del testo.
        \item \textbf{unsigned char *key}: La chiave AES 256 in formato binario.
        \item \textbf{unsigned char *iv}: L'IV del file.
    \end{itemize}
La funzione si occupa di inizzializzare il contesto di libssl e di cifrare la stringa plaintext, ovviamente la stringa viene cifrata mediante la chiave e l'IV. Dunque per ogni blocco la funzione si comprta nel seguente modo: procede a scrivere prima del blocco la dimensione di quest'ultimo, scrive il testo cifrato e infine scrive scrive il TAG di dimensione definita dalla costante TAG\_LEN, che in questa versione è grande esattamente 16 byte. Quest'ultimo ci garantisce l'integrità del blocco scritto poichè al suo interno conserva il valore del polinomio di Galois. Infine ritorna la dimensione del testo cifrato.
Di fatto ogni blocco difrato ha la seguente forma:
\begin{table}[htbp]
        \centering
        \begin{tabular}{|l|l|l|}
            \hline
            Lunghezza del testo cifrato (4 byte) & Testo cifrato & Tag (16 byte)\\
            \hline
        \end{tabular}
        \caption{Descrizione dei blocchi nella WRITE}
    \label{tab:tabella-blocchi-write}
\end{table}

\subsection{La decifratura dei blocchi}
La decifratura viene svolta dalla funzione decrypt\_file\_gcm che si aspetta i seguenti parametri:
 \begin{itemize}
        \item \textbf{FILE *fp}: Il puntatore al file da leggere.
        \item \textbf{unsigned char *key}: La chiave.
        \item \textbf{unsigned char *iv\_len}: L' initialization vector.
        \item \textbf{unsigned char **plaintext}: Un puntatore al buffer che deve contenere il testo una volta decifrato.
    \end{itemize}
La funzione si aspetta questa configurazione dei blocchi all'interno del file:
\begin{table}[htbp]
        \centering
        \begin{tabular}{|l|l|l|}
            \hline
            Lunghezza del testo cifrato (4 byte) & Testo cifrato & Tag (16 byte)\\
            \hline
        \end{tabular}
        \caption{Descrizione dei blocchi nella READ}
    \label{tab:tabella-blocchi-read}
\end{table}

Per decifrare il file la funzione legge la lunghezza del testo cifrato, successivamente legge il testo e il tag e lo decripta tramite la key e l'IV, dunque procede a scrivere il testo in \textit{*plaintext}. La funzione ritorna la lunghezza del testo in chiaro.

\section{L'implementazione della cifratura e della decifratura} 
TCFS Utilizza AES in modalità GCM per cifrare i blocchi dei file. Questo è una combinazione dell'algoritmo di cifratura simmetrica Advanced Encryption Standard (AES) con la modalità di autenticazione Galois/Counter Mode (GCM),  la modalità di funzionamento di questo algoritmo di cifratura è già stato adescritto nel capitolo precedente. Per completezza riportiamo adesso l'implementazione della lettura e della scrittura con AES GCM:
\begin{lstlisting}[style=C, caption=Esempio di decrypt e decrypt con gcm, label=C-gcm-expl]
extern int
decrypt_file_gcm (FILE *fp, unsigned char *key, unsigned char *iv,
                  unsigned char **plaintext, size_t bytes_to_read, off_t offset)
{
  EVP_CIPHER_CTX *ctx = NULL;
  int len;
  volatile int plaintext_len = 0;
  size_t aligned_offset = 0;
  size_t bytes_read = 0L;
  size_t bytes_read_prev = 0L;
  off_t seek_value = 0L;
  unsigned char *ciphertext = NULL;
  unsigned char tag[TAG_LEN];
  *plaintext = NULL;

  if (setjmp (jump_buffer))
    {
      logErr (
          "An error occurred in decrypt_file_gcm! Freeing resources... error");
      if (ctx)
        EVP_CIPHER_CTX_free (ctx);
      if (ciphertext)
        free (ciphertext);
      return false;
    }

  // Get file size
  fseek (fp, 0, SEEK_END);
  long fsize = ftell (fp);
  fseek (fp, 0, SEEK_SET); // same as rewind(f);

  // Align to offset
  if (offset > fsize)
    {
      aligned_offset = offset;
    }
  while (offset > aligned_offset && offset < fsize)
    {
      fread (&seek_value, sizeof (int), 1, fp);
      bytes_read += (unsigned long)sizeof (int);
      fseek (fp, seek_value, SEEK_CUR);
      bytes_read += seek_value;
      fseek (fp, TAG_LEN, SEEK_CUR);
      bytes_read += TAG_LEN;
      aligned_offset += seek_value;
      if (aligned_offset > offset)
        {
          logDebug ("Offset overshoot");
          bytes_read = bytes_read_prev;
          break;
        }
      else if (aligned_offset == offset)
        {
          logDebug ("Offset Found");
          break;
        }
      else
        {
          bytes_read_prev = bytes_read;
        }
    }

  fseek (fp, bytes_read, SEEK_SET);

  while (ftell (fp) < fsize && plaintext_len < bytes_to_read)
    {
      int ciphertext_len;
      fread (&ciphertext_len, sizeof (int), 1, fp);
      if (ciphertext_len < 0){
          handleErrors();
        }

      logDebug ("Allocating ciphertext of size %d", ciphertext_len);
      ciphertext = malloc (ciphertext_len);
      if (!ciphertext)
        {
          handleErrors ();
        }

      fread (ciphertext, sizeof (unsigned char), ciphertext_len, fp);

      fread (tag, sizeof (unsigned char), TAG_LEN,
             fp); // Read the tag from the file

      if (!(ctx = EVP_CIPHER_CTX_new ()))
        {
          handleErrors ();
        }

      if (1 != EVP_DecryptInit_ex (ctx, EVP_aes_256_gcm (), NULL, key, iv))
        {
          handleErrors ();
        }

      if (!EVP_CIPHER_CTX_ctrl (ctx, EVP_CTRL_GCM_SET_TAG, TAG_LEN, tag))
        {
          handleErrors ();
        }

      *plaintext = realloc (*plaintext, plaintext_len + ciphertext_len);
      if (!*plaintext)
        {
          handleErrors ();
        }

      if (1
          != EVP_DecryptUpdate (ctx, *plaintext + plaintext_len, &len,
                                ciphertext, ciphertext_len))
        {
          handleErrors ();
        }
      plaintext_len += len;

      EVP_CIPHER_CTX_free (ctx);

      free (ciphertext);
    }

  return plaintext_len;
}


  extern int
encrypt_file_gcm (FILE *fp, unsigned char *plaintext, size_t plaintext_len,
                  unsigned char *key, unsigned char *iv, off_t offset)
{
  EVP_CIPHER_CTX *ctx = NULL;
  int len;
  int ciphertext_len;
  unsigned char *ciphertext = NULL;
  unsigned char tag[TAG_LEN];

  if (setjmp (jump_buffer))
    {
      logErr (
          "An error occurred in encrypt_file_gcm! Freeing resources... error");
      EVP_CIPHER_CTX_free (ctx);
      if (ciphertext)
        free (ciphertext);
      return false;
    }

  // Align to offset
  off_t aligned_offset = (offset / TAG_LEN) * TAG_LEN;
  fseek (fp, aligned_offset, SEEK_SET);

  if (!(ctx = EVP_CIPHER_CTX_new ()))
    handleErrors ();

  if (1 != EVP_EncryptInit_ex (ctx, EVP_aes_256_gcm (), NULL, key, iv))
    handleErrors ();

  ciphertext = (unsigned char *)malloc (
      plaintext_len + EVP_CIPHER_block_size (EVP_aes_256_gcm ()));
  if (!ciphertext)
    handleErrors ();

  if (1 != EVP_EncryptUpdate (ctx, ciphertext, &len, plaintext, plaintext_len))
    handleErrors ();
  ciphertext_len = len;

  if (1 != EVP_EncryptFinal_ex (ctx, ciphertext + len, &len))
    handleErrors ();
  ciphertext_len += len;

  if (1 != EVP_CIPHER_CTX_ctrl (ctx, EVP_CTRL_GCM_GET_TAG, TAG_LEN, tag))
    handleErrors ();

  fwrite (&ciphertext_len, sizeof (int), 1, fp);

  if (fwrite (ciphertext, sizeof (unsigned char), ciphertext_len, fp)
      != (unsigned long)ciphertext_len)
    {
      handleErrors ();
    }
  if (fwrite (tag, sizeof (unsigned char), TAG_LEN, fp)
      != (unsigned long)TAG_LEN)
    {
      handleErrors ();
    }

  EVP_CIPHER_CTX_free (ctx);
  free (ciphertext);
  return ciphertext_len;
}
\end{lstlisting}
Ora spieghiamo il comportamento delle funzioni:
\subsubsection{decrypt\_file\_gcm}
La funzione decrypt\_file\_gcm è una funzione di decrittografia che utilizza l’algoritmo Galois/Counter Mode (GCM). Questa funzione prende come input un puntatore a un file, una chiave, un vettore di inizializzazione, un puntatore a un puntatore a un array di caratteri, il numero di byte da leggere e un offset. Inizialmente, la funzione inizializza un contesto di cifratura e altre variabili utilizzate per tenere traccia della posizione nel file e della lunghezza del testo in chiaro. Utilizza setjmp per gestire eventuali errori che possono verificarsi durante l’esecuzione della funzione. Se si verifica un errore, la funzione libera le risorse allocate e restituisce 0.
Successivamente, la funzione calcola la dimensione del file e si posiziona all’offset specificato mediante un ciclo. Questo viene fatto per calcolare il punto di salto nel file corrispondente all'offset specificato. Entra in un altro ciclo che continua fino a quando non ha letto tutti i byte richiesti o fino a quando non ha raggiunto la fine del file. In ogni iterazione del ciclo, legge la lunghezza del testo cifrato, il testo cifrato stesso e il TAG di autenticazione dal file.
Dopo aver letto il testo cifrato e il TAG, la funzione inizializza un nuovo contesto di cifratura e imposta la chiave e il vettore di inizializzazione. Imposta anche il TAG di autenticazione utilizzando EVP\_CIPHER\_CTX\_ctrl. Successivamente, decifra il testo cifrato e lo aggiunge al testo in chiaro. Se si verifica un errore durante la decifratura, la funzione gestisce l’errore e libera le risorse allocate.
Infine, alla fine del ciclo, la funzione libera il contesto di cifratura e il testo cifrato e restituisce la lunghezza del testo in chiaro decifrato.
\subsubsection{encrypt\_file\_gcm}
La funzione encrypt\_file\_gcm è utilizzata per crittografare un file utilizzando l’algoritmo di crittografia GCM (Galois/Counter Mode). La funzione inizia inizializzando un contesto di cifratura e alcune variabili. Utilizza setjmp per gestire eventuali errori che possono verificarsi durante l’esecuzione della funzione. Se si verifica un errore, la funzione libera le risorse allocate e restituisce false.
Successivamente, la funzione si posiziona all’offset specificato nel file. Crea un nuovo contesto di cifratura e imposta la chiave e il vettore di inizializzazione. Alloca memoria per il testo cifrato.
La funzione poi crittografa il testo in chiaro e lo aggiunge al testo cifrato. Se si verifica un errore durante la cifratura, la funzione gestisce l’errore e libera le risorse allocate.
Dopo aver cifrato il testo in chiaro, la funzione recupera il tag di autenticazione e scrive la lunghezza del testo cifrato, il testo cifrato e il tag nel file.
Infine, la funzione libera il contesto di cifratura e il testo cifrato e restituisce la lunghezza del testo cifrato. Questa funzione è un esempio di come si può utilizzare l’API OpenSSL per crittografare un file utilizzando l’algoritmo GCM.


\section{Uso di AES CTR per cifrare e decifrare i file}
L'utilizzo di AES CTR è molto simile a GCM con una singola differenza, non viene salvato il TAG dopo il blocco, di fatto i blocchi sul file hanno la seguente forma:\\
\begin{table}[htbp]
        \centering
        \begin{tabular}{|l|l|}
            \hline
            Lunghezza del testo cifrato (4 byte) & Testo cifrato \\
            \hline
        \end{tabular}
        \caption{Descrizione dei blocchi nella READ-WRITE di AES CTR}
    \label{tab:tabella-blocchi-read-write-ctr}
\end{table} \\
Le funzioni di lettura e scrittura sono molto simili alla modalità GCM (a meno del tag) e per questo non vengono riportate.

\subsection{La gerarchia e la generazione delle chiavi in TCFS}
Il FS di TCFS non utilizza la stessa chiave per cifrare ogni file. Quando un nuovo file viene creato e successivamente scritto, TCFS genera una nuova chiave di 32 byte, che viene cifrata tramite la chiave principale di TCFS ed inserita tra gli extended attributes del file. Questa chiave viene creata tramite la funzione generate\_key riportata di sotto:
\begin{lstlisting}[style=C, caption=Esempio di generate\_key, label=C-gen-key-expl]
void generate_key (unsigned char *destination)
{
  for (int i = 0; i < 10; i++)
    {
      int entropy = check_entropy ();
      if (entropy < 128)
        {
          logWarn ("Not enough entropy, creating some...");
          add_entropy ();
        }

      if (RAND_bytes (destination, 32) != true)
        {
          logErr ("Cannot generate key");
          destination = NULL;
        }

      if (strlen ((const char *)destination) == 32)
        break;
    }

  if (!is_valid_key (destination))
    {
      logErr ("ERR: Generated key is invalid\n");
      destination = NULL;
    }
}
\end{lstlisting}
La funzione check\_entropy() si occupa di verificare la presenza di un numero sufficiente di byte in "/dev/urandom" per la creazione della chiave, mentre add\_entropy() affronta un problema più specifico. In particolare, durante la creazione rapida di numerosi file, potrebbe verificarsi una situazione in cui "/dev/urandom" non sia in grado di generare nuovi byte con la stessa velocità. In questo contesto, la funzione add\_entropy() interviene attendendo che il sistema accumuli abbastanza entropia per garantire una generazione adeguata di byte.

\subsubsection{La differenza tra /dev/random e /dev/urandom}
Il dispositivo speciale "/dev/urandom" in ambienti Unix-like fornisce un flusso di dati pseudo-casuali utilizzati per scopi crittografici e di sicurezza.
A differenza di "/dev/random", che potrebbe bloccarsi se non c'è abbastanza entropia disponibile, /"dev/urandom" è progettato per essere un generatore di numeri casuali più rapido e non bloccante, tuttavia nel testing eseguito su TCFS si sono verificate due condizioni che hanno mandato in errore la generazione delle chiavi:
\begin{itemize}
    \item \textbf{La velocità di urandom non era sufficente}: Nonostante sia progettato per essere quanto più veloce possibile è stato notato che si poteva cadere nel caso in cui non erano stati generati abbastanza byte per ottenere una nuova chiave.
    \item \textbf{Non bloccante}: In caso si verifichi la condizione precedente allora il comportamento non bloccante forza urandom a restituire una sequenza più piccola di quella richiesta.
\end{itemize}
Dunque add\_entropy() interviene in situazioni in cui la generazione rapida di file potrebbe esaurire temporaneamente l'entropia disponibile, consentendo al sistema di accumulare entropia sufficiente per garantire una generazione affidabile dei byte necessari. Questo approccio garantisce una gestione efficace dell'entropia durante operazioni ad alta intensità di creazione di file.

\section{Uso di AES CBC per cifrare e decifrare i percorsi}
TCFS utilizza AES CBC per cifrare e decifrare i nomi delle cartelle. Abbiamo già discusso ampliamente di CBC dunque concentriamoci sulle funzioni che lo rendono possibile. Si utilizzano le seguenti funzioni della libreria crypt-utils:    
\begin{lstlisting}[style=C, caption=Funzioni di crypt-utils, label=C-crypt-utils.h]
const char *encrypt_path (const char *path, const char *key);

const char *encrypt_path_and_filename (const char *path, const char *key);

const char *decrypt_path (const char *encrypted_path, const char *key);

const char *decrypt_path_and_filename (const char *encrypted_path,
                                       const char *key);
\end{lstlisting}
chiaramente encrypt\_path e encrypt\_path\_and\_filename funzionano in maniera estamamente similare, la prima cripta il percorso fino alla cartella desiderata e la seconda cripta il percorso fino al file desiderato, lo stesso vale per decrypt\_path e decrypt\_path\_and\_filename.

\subsection{Perchè AES CBC si scontra con la scrittura dei file in FUSE}
Adesso spieghiamo perché si è deciso di non utilizzare CBC per cifrare il contenuto dei file:
AES CBC è un modo di operazione per la crittografia a blocchi che presenta alcune sfide quando viene utilizzato con FUSE. In particolare, il problema principale con l’uso di AES CBC in questo contesto è legato alla natura stessa di CBC. In CBC, ogni blocco di testo in chiaro viene combinato con il blocco di testo cifrato precedente prima di essere cifrato. Questo significa che la cifratura di un blocco dipende da tutti i blocchi precedenti. Pertanto, se si desidera modificare un singolo blocco all’interno di un file, non è possibile semplicemente decifrare e ricifrare quel blocco singolo. Invece, è necessario decifrare e ricifrare tutti i blocchi dal blocco modificato fino alla fine del file. Questo può essere molto inefficiente, soprattutto per file di grandi dimensioni.
Questo problema è particolarmente evidente quando si utilizza FUSE, che consente operazioni di lettura e scrittura a livello di file. Se un’applicazione tenta di scrivere in una posizione arbitraria all’interno di un file, con AES CBC l’intero file deve essere riscritto, il che può comportare un notevole sovraccarico. Inoltre la lettura randomica in AES CBC è semplicemente impossibile se non si decifra l'intero file.
Per questi motivi, AES CBC non è la scelta migliore per la crittografia a livello di directory o filesystem, ed è dunque stato considerato solo per la cifratura dei nomi delle cartelle e dei file.

\subsection{Dettagli implementativi di encrypt\_path\_and\_filename e decrypt\_path\_and\_filename}
Diamo ora un occhiata a encrypt\_path\_and\_filename per capire meglio la sua implementazione:
\begin{lstlisting}[style=C, caption=Funzione encrypt\_path\_and\_filename, label=C-enc-path-filename.h]
const char *encrypt_path_and_filename (const char *path, const char *key)
{
  char *result = NULL;
  char *token, *saveptr;

  // Check if the path is ".", ".."
  if (strcmp (path, ".") == 0 || strcmp (path, "..") == 0)
    {
      // If it is, no encryption is applied
      result = strdup (path);
      if (!result)
        {
          logErr ("Cannot allocate memory");
          exit (EXIT_FAILURE);
        }
      return result;
    }
  // Check if the path is /
  else if (strcmp (path, "/") == 0)
    {
      logDebug ("got root path\n");
      result = malloc (1 * sizeof (char));
      if (result == NULL)
        {
          logErr ("Cannot allocate memory");
          return NULL;
        }
      result[0] = '\0';
    }

  // Copy the original path
  char *path_copy = strdup (path);
  if (!path_copy)
    {
      logErr ("Cannot allocate memory");
      exit (EXIT_FAILURE);
    }

  // Start encryption
  token = strtok_r (path_copy, "/", &saveptr);
  while (token != NULL)
    {
      // Check if the directory is "." or ".."
      if (strcmp (token, ".") == 0 || strcmp (token, "..") == 0)
        {
          // If it is, no encryption is applied
        }
      else
        {
          // Encrypt each part of the path
          const char *encrypted_part = encrypt_file_name_with_hex (token, key);
          logDebug ("\tEncrypted %s --> %s\n", token, encrypted_part);

          // Concatenate to the result string
          if (result == NULL)
            {
              // Check if the first character of encrypted_part is '/'
              if (encrypted_part[0] != '/')
                {
                  result = malloc (strlen (encrypted_part) + 2);
                  if (!result)
                    {
                      logErr ("Cannot allocate memory");
                      exit (EXIT_FAILURE);
                    }
                  strcpy (result, "/");
                  strcat (result, encrypted_part);
                }
              else
                {
                  result = strdup (encrypted_part);
                }
            }
          else
            {
              size_t result_len = strlen (result);
              size_t encrypted_len = strlen (encrypted_part);

              result = realloc (result, result_len + 1 + encrypted_len + 1);
              if (!result)
                {
                  logErr ("Cannot allocate memory");
                  exit (EXIT_FAILURE);
                }

              // Check if the first character of encrypted_part is '/'
              if (encrypted_part[0] != '/')
                {
                  strcat (result, "/");
                }

              strcat (result, encrypted_part);
              logDebug ("Tempresult: %s\n", result);
            }
        }

      // Move to the next part of the path
      token = strtok_r (NULL, "/", &saveptr);
    }

  // Free the memory allocated for the path copy
  free (path_copy);

  return result;
}
\end{lstlisting}

Notiamo innanzitutto che ci sono casi speciali, essa controlla se il percorso del file è uno dei casi speciali “.”, “..” o “/”. In questi casi, non applica nessuna criptazione e restituisce il percorso originale come risultato, dopo averlo copiato con la funzione strdup.
Se il percorso del file non è uno dei casi speciali, la funzione lo copia in una variabile locale 
path\_copy, usando ancora la funzione strdup.
La funzione inizia il processo di criptazione, usando la funzione strtok\_r per dividere il percorso del file in base al carattere “/”. Questo significa che la funzione considera ogni parte del percorso come una directory da criptare separatamente. Per ogni parte del percorso, la funzione controlla se si tratta di una delle directory speciali “.” o “..”. In questo caso, non applica nessuna criptazione e passa alla parte successiva.
Se la parte del percorso non è una directory speciale, la funzione la cripta usando una funzione ausiliaria chiamata encrypt\_file\_name\_with\_hex, che prende in input la parte del percorso e la chiave di criptazione. Questa funzione restituisce una stringa che rappresenta la parte criptata in esadecimale.
La funzione concatena la parte criptata alla stringa result, aggiungendo il carattere “/” se necessario. Se result è ancora NULL, significa che è la prima parte del percorso da criptare, e quindi la funzione la inizializza con la parte criptata. Altrimenti, la funzione usa la funzione realloc per aumentare la dimensione di result e aggiungere la parte criptata alla fine.
La funzione ripete il processo finché non ha esaurito tutte le parti del percorso del file, usando la variabile token per memorizzare la parte corrente e la variabile saveptr per memorizzare il resto del percorso da processare. La funzione libera la memoria allocata per la copia del percorso del file, usando la funzione free e restituisce la stringa result, che contiene il percorso e il nome del file criptati.
La funzione decrypt\_path\_and\_filename funziona esattamente allo stesso modo ma invece di invocare la funzione ausiliaria encrypt\_file\_name\_with\_hex viene invocata la decrypt\_file\_name\_with\_hex.

\subsection{Le funzioni ausiliarie encrypt\_file\_name\_with\_hex e decrypt\_file\_name\_with\_hex}
Queste due funzioni fungono da facade per le funzioni encrypt\_string e decrypt\_string. analizziamole semparatamente.

\subsubsection{La funzione encrypt\_string}
Questa funzione ha lo scopo di inizializzare il contesto di AES CBC, criptare una stringa di lunghezza arbitraria e restituisce la stringa criptata in esadecimale. Diamo anche qui un occhiata alla sua implementazione:
\begin{lstlisting}[style=C, caption=Funzione encrypt\_string, label=C-key-utils.h]
unsigned char *encrypt_string (unsigned char *plaintext, const char *key,
                int *encrypted_key_len)
{
  EVP_CIPHER_CTX *ctx;
  const EVP_CIPHER *cipher = EVP_aes_256_cbc ();
  unsigned char iv[AES_BLOCK_SIZE];
  memset (iv, 0, AES_BLOCK_SIZE);

  ctx = EVP_CIPHER_CTX_new ();
  if (!ctx)
    {
      return NULL;
    }

  EVP_EncryptInit_ex (ctx, cipher, NULL, (const unsigned char *)key, iv);

  size_t plaintext_len = strlen ((const char *)plaintext);
  unsigned char ciphertext[plaintext_len + AES_BLOCK_SIZE];
  memset (ciphertext, 0, sizeof (ciphertext));

  int len, total_len = 0;
  EVP_EncryptUpdate (ctx, ciphertext, &len, plaintext, (int)plaintext_len);
  total_len += len;
  EVP_EncryptFinal_ex (ctx, ciphertext + total_len, &len);
  total_len += len;
  EVP_CIPHER_CTX_free (ctx);

  unsigned char *encoded_string = malloc (total_len * 2 + 1);
  if (!encoded_string)
    {
      logErr ("Cannot allocate memory for encrypt_string encoded_string");
      return NULL;
    }

  for (int i = 0; i < total_len; i++)
    {
      sprintf ((char *)&encoded_string[i * 2], "%02x", ciphertext[i]);
    }
  encoded_string[total_len * 2] = '\0';

  *encrypted_key_len = total_len * 2;
  return encoded_string;
}
\end{lstlisting}
Per iniziare, la funzione crea una struttura che contiene le informazioni necessarie per la criptazione, come il tipo di metodo AES, la parola chiave e il valore casuale. Questa struttura si chiama contesto di criptazione e viene inizializzata con le funzioni apposite. Successivamente la funzione calcola la lunghezza della frase da criptare e crea un vettore di byte che ha abbastanza spazio per contenere il codice segreto. Dopo, la funzione inizia a criptare la frase e a scrivere il codice segreto nel vettore di byte, usando le funzioni di criptazione. Queste funzioni prendono il contesto di criptazione, la frase e il vettore di byte e restituiscono il numero di byte scritti. La funzione tiene conto di questo numero per sapere quanto è lungo il codice segreto.
Infine, la funzione libera la memoria occupata dal contesto di criptazione e converte il codice segreto in una stringa di caratteri esadecimali. Questo significa che ogni byte viene rappresentato da due cifre da 0 a 9 e da A a F. La funzione restituisce questa stringa, che è il risultato della criptazione.

\subsubsection{La funzione decrypt\_string}

\begin{lstlisting}[style=C, caption=Funzione decrypt\_string, label=C-dec-string.h]
unsigned char *
decrypt_string (unsigned char *ciphertext, const char *key)
{
  EVP_CIPHER_CTX *ctx;
  const EVP_CIPHER *cipher = EVP_aes_256_cbc ();
  unsigned char iv[AES_BLOCK_SIZE];
  memset (iv, 0, AES_BLOCK_SIZE);

  ctx = EVP_CIPHER_CTX_new ();
  EVP_DecryptInit_ex (ctx, cipher, NULL, (const unsigned char *)key, iv);

  size_t ciphertext_len = strlen ((const char *)ciphertext) / 2;
  unsigned char *decoded_ciphertext = malloc (ciphertext_len);

  if (!decoded_ciphertext)
    {
      logErr ("Cannot allocate memory for decrypt_string decoded_ciphertext");
      return NULL;
    }

  for (size_t i = 0; i < ciphertext_len; i++)
    {
      char hex[3]
          = { (char)ciphertext[i * 2], (char)ciphertext[i * 2 + 1], '\0' };
      char *endptr;
      unsigned long byte = strtoul (hex, &endptr, 16);

      if (*endptr != '\0' || byte > UCHAR_MAX)
        {
          logErr ("decrypt string error");
          free ((void *)decoded_ciphertext);
          return NULL;
        }
      decoded_ciphertext[i] = byte;
    }

  unsigned char plaintext[ciphertext_len + AES_BLOCK_SIZE + 1];
  memset (plaintext, 0, sizeof (plaintext));

  int len;
  EVP_DecryptUpdate (ctx, plaintext, &len, decoded_ciphertext,
                     (int)ciphertext_len);
  int padding_len;
  EVP_DecryptFinal_ex (ctx, plaintext + len, &padding_len);
  EVP_CIPHER_CTX_free (ctx);

  unsigned char *decrypted_string
      = (unsigned char *)malloc (len + padding_len + 1);

  if (!decrypted_string)
    {
      logErr ("Cannot allocate memory for decrypt_string decrypted_string");
      return NULL;
    }

  memcpy (decrypted_string, plaintext, len + padding_len);
  decrypted_string[len + padding_len] = '\0';

  free (decoded_ciphertext);

  return decrypted_string;
}
\end{lstlisting}
Per decriptare il codice segreto, la funzione deve prima convertirlo in un vettore di byte, cioè una sequenza di byte che può essere manipolata da libssl. Per fare questo, la funzione usa un ciclo for che scorre la stringa esadecimale e prende ogni coppia di caratteri. Poi, usa la funzione strtoul per convertire la coppia di caratteri in un valore numerico, interpretando la stringa come un numero esadecimale. Se la conversione ha successo, la funzione assegna il valore numerico al byte corrispondente nel vettore di byte. Se la conversione fallisce o il valore numerico è troppo grande per essere contenuto in un byte, la funzione restituisce un errore e termina.
Dopo aver convertito il codice segreto in un vettore di byte, la funzione deve creare una struttura che contiene le informazioni necessarie per la decrittazione, come il tipo di metodo AES, la parola chiave e il valore casuale. Questa struttura si chiama contesto di decrittazione e viene usata dalle funzioni di decrittazione per eseguire le operazioni inverse di quelle fatte dalla criptazione. Il tipo di metodo AES è una costante che indica il tipo di cifrario AES da usare (in questo caso AES-256-CBC, che significa che usa una chiave di 256 bit e una modalità di operazione a blocchi concatenati). La parola chiave è una stringa che rappresenta la chiave segreta di decrittazione, che deve essere la stessa usata per la criptazione. Il valore casuale è un vettore di byte che rappresenta il vettore di inizializzazione, che è un valore usato per rendere più sicura la decrittazione. In questo caso, il vettore di inizializzazione è inizializzato con tutti zeri, usando la funzione memset.
Per inizializzare il contesto di decrittazione, la funzione usa due funzioni apposite: EVP\_CIPHER\_CTX\_new e EVP\_DecryptInit\_ex. La prima funzione crea un nuovo contesto di decrittazione e restituisce un puntatore alla struttura. La seconda funzione inizializza il contesto di decrittazione con il tipo di metodo AES, la parola chiave e il vettore di inizializzazione.
Dopo aver inizializzato il contesto di decrittazione, la funzione deve creare un vettore di byte che ha abbastanza spazio per contenere la frase decrittata. Il vettore di byte ha una dimensione pari alla lunghezza del codice segreto diviso due (perché ogni byte corrisponde a due caratteri esadecimali) più la dimensione di un blocco AES (16 byte), per garantire che ci sia abbastanza spazio per contenere la frase decrittata. Il vettore di byte viene riempito con dei valori nulli, usando la funzione memset, per evitare che ci siano dei dati inutili.
Poi, la funzione inizia a decrittare il codice segreto e a scrivere la frase nel vettore di byte, usando le funzioni di decrittazione. Queste funzioni sono EVP\_DecryptUpdate e EVP\_DecryptFinal\_ex. La prima funzione decritta una parte del codice segreto e la scrive nel vettore di byte, restituendo il numero di byte scritti. La seconda funzione decritta la parte finale del codice segreto e la scrive nel vettore di byte, restituendo il numero di byte scritti. La funzione tiene conto di questi numeri per sapere quanto è lunga la frase decrittata.
Infine, la funzione libera la memoria occupata dal contesto di decrittazione, usando la funzione EVP\_CIPHER\_CTX\_free, e copia la frase dal vettore di byte in una stringa di caratteri, usando la funzione malloc per allocare la memoria e la funzione memcpy per copiare i dati. Questa stringa viene terminata con il carattere '\textbackslash 0', che indica la fine della frase. La funzione restituisce questa stringa, che è il risultato della decrittazione.